\chapter{Code Source}

\section{Algorithme A* pour Plus Court Chemin}

L'implémentation de l'algorithme A* avec fonction de coût composite est présentée ci-dessous. Le code utilise une file de priorité pour explorer les nœuds par ordre croissant de score $f(n) = g(n) + h(n)$.

\begin{lstlisting}[language=Java, caption={Algorithme A* - AStar.java}, label={lst:astar}]
package com.delivery.optimization.algorithm;

import com.delivery.optimization.domain.Arc;
import com.delivery.optimization.domain.Node;
import lombok.Value;
import org.springframework.stereotype.Component;
import java.util.*;

@Component
public class AStar {
    @Value
    public static class PathResult {
        List<String> path;
        double totalCost;
    }

    public PathResult findPath(String originId, String destinationId,
            Map<String, Node> nodes, Map<String, List<Arc>> adjacencyList,
            CostFunction.Weights weights) {
        PriorityQueue<NodeScore> openSet = new PriorityQueue<>(
            Comparator.comparingDouble(NodeScore::getFScore));
        Map<String, String> cameFrom = new HashMap<>();
        Map<String, Double> gScore = new HashMap<>();

        gScore.put(originId, 0.0);
        openSet.add(new NodeScore(originId,
            heuristic(originId, destinationId, nodes)));

        while (!openSet.isEmpty()) {
            NodeScore current = openSet.poll();
            String currentId = current.getNodeId();

            if (currentId.equals(destinationId)) {
                return new PathResult(reconstructPath(cameFrom, currentId),
                    gScore.get(currentId));
            }

            for (Arc arc : adjacencyList.getOrDefault(currentId,
                    Collections.emptyList())) {
                String neighborId = arc.getDestinationId();
                double tentativeGScore = gScore.get(currentId)
                    + calculateCost(arc, weights);

                if (tentativeGScore < gScore.getOrDefault(neighborId,
                        Double.MAX_VALUE)) {
                    cameFrom.put(neighborId, currentId);
                    gScore.put(neighborId, tentativeGScore);
                    double fScore = tentativeGScore
                        + heuristic(neighborId, destinationId, nodes);

                    if (openSet.stream().noneMatch(
                            ns -> ns.getNodeId().equals(neighborId))) {
                        openSet.add(new NodeScore(neighborId, fScore));
                    }
                }
            }
        }
        return null; // Path not found
    }

    private double calculateCost(Arc arc, CostFunction.Weights weights) {
        double tf = arc.getTrafficFactor() != null
            ? arc.getTrafficFactor() : 1.0;
        return weights.getAlpha() * arc.getDistance()
            + weights.getBeta() * (arc.getTravelTime() * tf)
            + weights.getGamma() * arc.getPenibility()
            + weights.getDelta() * arc.getWeatherImpact()
            + weights.getEta() * arc.getFuelCost();
    }

    private double heuristic(String nodeId, String targetId,
            Map<String, Node> nodes) {
        Node n = nodes.get(nodeId);
        Node d = nodes.get(targetId);
        if (n == null || d == null) return 0;

        double dHaversine = haversine(n.getLatitude(), n.getLongitude(),
            d.getLatitude(), d.getLongitude());
        double vmax = 60.0; // Max speed in km/h
        return dHaversine / vmax;
    }

    private double haversine(double lat1, double lon1,
            double lat2, double lon2) {
        double R = 6371; // Earth radius in km
        double dLat = Math.toRadians(lat2 - lat1);
        double dLon = Math.toRadians(lon2 - lon1);
        double a = Math.sin(dLat / 2) * Math.sin(dLat / 2) +
            Math.cos(Math.toRadians(lat1)) * Math.cos(Math.toRadians(lat2)) *
            Math.sin(dLon / 2) * Math.sin(dLon / 2);
        double c = 2 * Math.atan2(Math.sqrt(a), Math.sqrt(1 - a));
        return R * c;
    }

    private List<String> reconstructPath(Map<String, String> cameFrom,
            String current) {
        List<String> path = new ArrayList<>();
        path.add(current);
        while (cameFrom.containsKey(current)) {
            current = cameFrom.get(current);
            path.add(0, current);
        }
        return path;
    }
}
\end{lstlisting}

\section{Filtre de Kalman pour Prédiction ETA}

L'implémentation du filtre de Kalman étendu avec biais adaptatif utilise Apache Commons Math pour les opérations matricielles.

\begin{lstlisting}[language=Java, caption={Filtre de Kalman - KalmanFilter.java}, label={lst:kalman}]
package com.delivery.optimization.algorithm;

import lombok.Builder;
import lombok.Data;
import org.apache.commons.math3.linear.*;
import org.springframework.stereotype.Component;

@Component
public class KalmanFilter {
    @Data
    @Builder
    public static class ExtendedState {
        private RealVector x; // [distance, speed, bias]
        private RealMatrix p; // covariance matrix
    }

    private static final double ALPHA_B = 0.95; // Bias decay factor

    /**
     * Prediction step: x_k|k-1 = F * x_k-1, P_k|k-1 = F * P_k-1 * F^T + Q
     */
    public ExtendedState predict(ExtendedState state, double dt,
            RealMatrix q) {
        // F matrix as defined in Section 3.6.2
        RealMatrix f = new Array2DRowRealMatrix(new double[][] {
            { 1.0, dt, 0.0 },
            { 0.0, 1.0, 0.0 },
            { 0.0, 0.0, ALPHA_B }
        });

        RealVector predictedX = f.operate(state.getX());
        RealMatrix predictedP = f.multiply(state.getP())
            .multiply(f.transpose()).add(q);

        return ExtendedState.builder()
            .x(predictedX)
            .p(predictedP)
            .build();
    }

    /**
     * Update step: x_k|k = x_k|k-1 + K(z - H*x_k|k-1)
     */
    public ExtendedState update(ExtendedState state, double distMeasured,
            double speedMeasured, double rDist, double rSpeed) {
        // H matrix: Observe distance and speed
        RealMatrix h = new Array2DRowRealMatrix(new double[][] {
            { 1.0, 0.0, 0.0 },
            { 0.0, 1.0, 0.0 }
        });

        RealVector z = new ArrayRealVector(
            new double[] { distMeasured, speedMeasured });
        RealMatrix r = new Array2DRowRealMatrix(new double[][] {
            { rDist, 0.0 },
            { 0.0, rSpeed }
        });

        // Innovation y = z - Hx
        RealVector y = z.subtract(h.operate(state.getX()));

        // S = HPH^T + R
        RealMatrix s = h.multiply(state.getP())
            .multiply(h.transpose()).add(r);

        // K = PH^T * S^-1
        RealMatrix k = state.getP().multiply(h.transpose())
            .multiply(new LUDecomposition(s).getSolver().getInverse());

        // Updated state x = x + Ky
        RealVector updatedX = state.getX().add(k.operate(y));

        // Updated covariance P = (I - KH) * P
        RealMatrix i = MatrixUtils.createRealIdentityMatrix(3);
        RealMatrix updatedP = i.subtract(k.multiply(h))
            .multiply(state.getP());

        return ExtendedState.builder()
            .x(updatedX)
            .p(updatedP)
            .build();
    }
}
\end{lstlisting}

\newpage
\section{Solveur VRP avec Google OR-Tools}

Extrait de l'implémentation du solveur VRP utilisant Google OR-Tools pour l'optimisation de tournées avec contraintes de capacité et fenêtres temporelles.

\begin{lstlisting}[language=Java, caption={Solveur VRP - VRPSolver.java (extrait)}, label={lst:vrp}]
package com.delivery.optimization.algorithm;

import com.google.ortools.Loader;
import com.google.ortools.constraintsolver.*;
import lombok.extern.slf4j.Slf4j;
import org.springframework.stereotype.Component;

@Component
@Slf4j
public class VRPSolver {
    static {
        Loader.loadNativeLibraries(); // Charger OR-Tools natif
    }

    public TourOptimizationResponse solve(TourOptimizationRequest request,
            List<Node> availableRelays, Map<String, Node> allNodes,
            List<Arc> allArcs) {
        try {
            // 1. Preparer les donnees pour OR-Tools
            DataModel data = createDataModel(request, availableRelays,
                allNodes, allArcs);

            // 2. Creer le gestionnaire de routing
            RoutingIndexManager manager = new RoutingIndexManager(
                data.distanceMatrix.length, 1, data.depot);

            // 3. Creer le modele de routing
            RoutingModel routing = new RoutingModel(manager);

            // 4. Definir la fonction de cout (distance/temps)
            final int transitCallbackIndex = routing.registerTransitCallback(
                (long fromIndex, long toIndex) -> {
                    int fromNode = manager.indexToNode(fromIndex);
                    int toNode = manager.indexToNode(toIndex);
                    return data.distanceMatrix[fromNode][toNode];
                });
            routing.setArcCostEvaluatorOfAllVehicles(transitCallbackIndex);

            // 5. Ajouter contrainte de capacite
            final int demandCallbackIndex = routing.registerUnaryTransitCallback(
                (long fromIndex) -> {
                    int fromNode = manager.indexToNode(fromIndex);
                    return data.demands[fromNode];
                });
            routing.addDimensionWithVehicleCapacity(
                demandCallbackIndex, 0,
                new long[]{request.getVehicleCapacity()},
                true, "Capacity");

            // 6. Parametres de recherche
            RoutingSearchParameters searchParameters =
                main.defaultRoutingSearchParameters().toBuilder()
                .setFirstSolutionStrategy(
                    FirstSolutionStrategy.Value.PATH_CHEAPEST_ARC)
                .setLocalSearchMetaheuristic(
                    LocalSearchMetaheuristic.Value.GUIDED_LOCAL_SEARCH)
                .setTimeLimit(com.google.protobuf.Duration.newBuilder()
                    .setSeconds(30).build())
                .build();

            // 7. Resoudre
            Assignment solution = routing.solveWithParameters(searchParameters);

            if (solution == null) {
                log.warn("Aucune solution trouvee pour le VRP");
                return createFallbackSolution(request, availableRelays);
            }

            // 8. Extraire la solution
            return extractSolution(solution, routing, manager, data, request);

        } catch (Exception e) {
            log.error("Erreur lors de la resolution VRP: {}",
                e.getMessage(), e);
            return createFallbackSolution(request, availableRelays);
        }
    }
}
\end{lstlisting}

\newpage
\section{Scripts de Base de Données}

\subsection{Création de la Table Nodes}

\begin{lstlisting}[language=SQL, caption={Migration 001 - Création table nodes}, label={lst:sql-nodes}]
CREATE TABLE nodes (
    id VARCHAR(50) PRIMARY KEY,
    type VARCHAR(20) NOT NULL, -- CLIENT, RELAY, DEPOT
    name VARCHAR(100),
    latitude DOUBLE PRECISION NOT NULL,
    longitude DOUBLE PRECISION NOT NULL,
    capacity INT DEFAULT 0,
    current_occupancy INT DEFAULT 0
);

CREATE INDEX idx_nodes_type ON nodes(type);
\end{lstlisting}

\subsection{Tables Auxiliaires pour Métriques}

\begin{lstlisting}[language=SQL, caption={Migration 006 - Tables métriques et événements}, label={lst:sql-metrics}]
CREATE TABLE relay_points (
    id VARCHAR(50) PRIMARY KEY,
    name VARCHAR(100),
    latitude DOUBLE PRECISION NOT NULL,
    longitude DOUBLE PRECISION NOT NULL,
    capacity INT DEFAULT 0,
    current_occupancy INT DEFAULT 0
);

CREATE TABLE eta_history (
    id SERIAL PRIMARY KEY,
    delivery_id VARCHAR(50) REFERENCES deliveries(id),
    estimated_eta TIMESTAMP WITH TIME ZONE,
    actual_arrival TIMESTAMP WITH TIME ZONE,
    error_margin DOUBLE PRECISION,
    timestamp TIMESTAMP WITH TIME ZONE DEFAULT CURRENT_TIMESTAMP
);

CREATE TABLE reroute_events (
    id SERIAL PRIMARY KEY,
    delivery_id VARCHAR(50) REFERENCES deliveries(id),
    reason VARCHAR(255),
    old_path_cost DOUBLE PRECISION,
    new_path_cost DOUBLE PRECISION,
    hysteresis_met BOOLEAN,
    timestamp TIMESTAMP WITH TIME ZONE DEFAULT CURRENT_TIMESTAMP
);
\end{lstlisting}

\subsection{Migration Idempotente pour Filtre de Kalman}

\begin{lstlisting}[language=SQL, caption={Migration 011 - Ajout colonne total\_distance}, label={lst:sql-kalman}]
-- Add total_distance column if it doesn't exist
DO $$
BEGIN
    IF NOT EXISTS (
        SELECT 1 FROM information_schema.columns
        WHERE table_name='kalman_states'
        AND column_name='total_distance'
    ) THEN
        ALTER TABLE kalman_states
        ADD COLUMN total_distance DOUBLE PRECISION;
    END IF;
END $$;
\end{lstlisting}

\newpage
\chapter{Données Brutes et Résultats Détaillés}

\section{Structure des Tables PostgreSQL}

\subsection{Base delivery\_db}

La base de données \texttt{delivery\_db} contient les tables suivantes :

\begin{table}[h]
\centering
\caption{Tables principales de delivery\_db}
\label{tab:db-schema}
\small
\begin{tabular}{|l|p{10cm}|}
\hline
\textbf{Table} & \textbf{Description} \\
\hline
\texttt{nodes} & 67 nœuds du réseau (54 clients, 12 relais, 1 dépôt) avec coordonnées GPS \\
\hline
\texttt{arcs} & ~200 segments routiers avec distance, temps, pénibilité, météo, carburant \\
\hline
\texttt{deliveries} & Livraisons avec pickup/dropoff, poids, statut, horodatages \\
\hline
\texttt{drivers} & Coursiers avec position GPS actuelle, statut disponibilité \\
\hline
\texttt{tours} & Tournées optimisées avec coût total, distance, durée estimée \\
\hline
\texttt{tour\_stops} & Arrêts de tournée avec ordre de séquence, type (PICKUP/DROPOFF) \\
\hline
\texttt{kalman\_states} & États filtre de Kalman (distance, vitesse, biais, covariance) \\
\hline
\texttt{eta\_history} & Historique ETA estimés vs arrivées réelles pour calibration \\
\hline
\texttt{reroute\_events} & Événements de reroutage avec coûts ancien/nouveau chemins \\
\hline
\texttt{relay\_points} & Points relais avec capacité et occupation actuelle \\
\hline
\end{tabular}
\end{table}

\subsection{Base petri\_db}

La base de données \texttt{petri\_db} contient les tables pour modélisation des Réseaux de Petri :

\begin{table}[h]
\centering
\caption{Tables de petri\_db}
\label{tab:petri-schema}
\begin{tabular}{|l|p{10cm}|}
\hline
\textbf{Table} & \textbf{Description} \\
\hline
\texttt{petri\_nets} & Définitions de réseaux de Petri (workflow livraison) \\
\hline
\texttt{places} & Places du réseau (états : PENDING, ASSIGNED, IN\_TRANSIT, etc.) \\
\hline
\texttt{transitions} & Transitions entre places (événements métier) \\
\hline
\texttt{tokens} & Tokens actifs référençant les livraisons (entity\_id) \\
\hline
\texttt{arcs\_petri} & Arcs du réseau de Petri (place → transition, transition → place) \\
\hline
\end{tabular}
\end{table}

\section{Exemples de Données de Test}

\subsection{Nœuds du Réseau de Yaoundé}

Extrait du fichier \texttt{009-complete-seed.sql} montrant des nœuds réels de Yaoundé :

\begin{lstlisting}[language=SQL, caption={Exemples de nœuds clients}, label={lst:data-nodes}]
INSERT INTO nodes (id, type, name, latitude, longitude) VALUES
('depot_1', 'DEPOT', 'Depot Central Yaounde', 3.8480, 11.5021),
('client_1', 'CLIENT', 'Quartier Bastos', 3.8850, 11.5180),
('client_2', 'CLIENT', 'Quartier Melen', 3.8277, 11.5173),
('client_3', 'CLIENT', 'Quartier Mvog-Ada', 3.8500, 11.5100),
('relay_1', 'RELAY', 'Relais Mokolo Market', 3.8650, 11.5050, 50, 0),
('relay_2', 'RELAY', 'Relais Essos', 3.8967, 11.5119, 30, 0)
ON CONFLICT (id) DO NOTHING;
\end{lstlisting}

\subsection{Configuration Application Spring Boot}

Fichier \texttt{application.yml} montrant l'utilisation des variables Railway :

\begin{lstlisting}[language=yaml, caption={Configuration application.yml}, label={lst:config-yml}]
spring:
  application:
    name: delivery-optimization-api
  r2dbc:
    url: r2dbc:postgresql://${PGHOST:localhost}:${PGPORT:5432}/${PGDATABASE:delivery_db}
    username: ${PGUSER:postgres}
    password: ${PGPASSWORD:postgres}
    pool:
      initial-size: 5
      max-size: 20
  liquibase:
    enabled: true
    url: jdbc:postgresql://${PGHOST:localhost}:${PGPORT:5432}/${PGDATABASE:delivery_db}
    user: ${PGUSER:postgres}
    password: ${PGPASSWORD:postgres}
    change-log: classpath:db/changelog/db.changelog-master.xml

management:
  endpoints:
    web:
      exposure:
        include: health,info,prometheus,metrics

petri-net:
  api:
    url: ${PETRI_NET_API_URL:http://localhost:8081}
\end{lstlisting}

\section{Versions Exactes des Technologies}

\subsection{Frontend (package.json)}

\begin{lstlisting}[language=json, caption={Versions Frontend - package.json (extrait)}, label={lst:package-json}]
{
  "name": "delivery-optimization-frontend",
  "version": "0.1.0",
  "dependencies": {
    "@stomp/stompjs": "^7.2.1",
    "@tanstack/react-query": "^5.17.0",
    "leaflet": "^1.9.4",
    "next": "^14.1.0",
    "react": "^18.2.0",
    "react-dom": "^18.2.0",
    "react-leaflet": "^4.2.1",
    "recharts": "^2.10.0",
    "sockjs-client": "^1.6.1",
    "tailwindcss": "^3.4.0",
    "zustand": "^4.5.0"
  },
  "devDependencies": {
    "typescript": "^5.3.3"
  }
}
\end{lstlisting}

\subsection{Backend Delivery API (pom.xml)}

\begin{lstlisting}[language=xml, caption={Versions Backend Delivery - pom.xml (extrait)}, label={lst:pom-delivery}]
<parent>
    <groupId>org.springframework.boot</groupId>
    <artifactId>spring-boot-starter-parent</artifactId>
    <version>3.2.2</version>
</parent>

<properties>
    <java.version>17</java.version>
    <mapstruct.version>1.5.5.Final</mapstruct.version>
</properties>

<dependencies>
    <dependency>
        <groupId>org.springframework.boot</groupId>
        <artifactId>spring-boot-starter-webflux</artifactId>
    </dependency>
    <dependency>
        <groupId>org.apache.commons</groupId>
        <artifactId>commons-math3</artifactId>
        <version>3.6.1</version>
    </dependency>
    <dependency>
        <groupId>com.google.ortools</groupId>
        <artifactId>ortools-java</artifactId>
        <version>9.8.3296</version>
    </dependency>
</dependencies>
\end{lstlisting}

\subsection{Backend Petri Net API (pom.xml)}

\begin{lstlisting}[language=xml, caption={Versions Backend Petri Net - pom.xml (extrait)}, label={lst:pom-petri}]
<parent>
    <groupId>org.springframework.boot</groupId>
    <artifactId>spring-boot-starter-parent</artifactId>
    <version>3.4.0</version>
</parent>

<properties>
    <java.version>17</java.version>
</properties>

<dependencies>
    <dependency>
        <groupId>org.springframework.boot</groupId>
        <artifactId>spring-boot-starter-webflux</artifactId>
    </dependency>
    <dependency>
        <groupId>org.springframework.boot</groupId>
        <artifactId>spring-boot-starter-data-r2dbc</artifactId>
    </dependency>
    <dependency>
        <groupId>org.postgresql</groupId>
        <artifactId>r2dbc-postgresql</artifactId>
    </dependency>
</dependencies>
\end{lstlisting}